% \section{Sobre a CC2020 da ACM}

% O \textit{Computing Curricula 2020} (CC2020) da ACM (\textit{Association for Computing Machinery}), em conjunto com outras entidades como a IEEE (\textit{Institute of Electrical and Electronics Engineers}), estabelece diretrizes para a educação em computação, diferenciando claramente os diversos campos dentro da área. Ele fornece um modelo unificado que identifica conhecimentos essenciais, habilidades e competências, garantindo que a formação acadêmica esteja alinhada às demandas do mercado e aos avanços tecnológicos.

% Uma das principais contribuições do CC2020 é a diferenciação entre os diversos cursos na área da computação, ajudando a estruturar currículos de maneira específica para cada especialidade. As principais áreas abordadas pelo CC2020 são:

%     \begin{itemize}
%         \item \textbf{\textit{Computer Engineering} (Engenharia de Computação)}: Integra hardware e software, abordando sistemas embarcados, redes, eletrônica digital, computação paralela e arquiteturas computacionais.
%         \item \textbf{\textit{Computer Science} (Ciência da Computação)}: Foca nos fundamentos teóricos e práticos da computação, incluindo algoritmos, complexidade computacional, linguagens de programação e paradigmas computacionais.
%         \item \textbf{\textit{Cybersecurity} (Cibersegurança)}: Explora a proteção de sistemas, redes e dados contra ameaças e ataques, abordando criptografia, segurança da informação, defesa cibernética e gestão de riscos.
%         \item \textbf{\textit{Information Systems} (Sistemas de Informação)}: Analisa a relação entre computação e organizações, focando na gestão de dados, processos de negócios e suporte à tomada de decisão.
%         \item \textbf{\textit{Information Technology} (Tecnologia da Informação)}: Concentra-se na infraestrutura de TI, redes, suporte técnico, administração de sistemas e segurança operacional.
%         \item \textbf{\textit{Software Engineering} (Engenharia de Software)}: Foca no desenvolvimento estruturado de software, abrangendo metodologias de desenvolvimento, qualidade, testes e escalabilidade de sistemas.
%         \item \textbf{\textit{Data Science} (Ciência de Dados)}: Ainda em evolução dentro do CC2020; explora o uso de estatística, aprendizado de máquina e engenharia de dados para extrair insights de grandes volumes de dados.
%     \end{itemize}

% Para um curso de Engenharia de Computação que segue as Diretrizes Curriculares Nacionais de 2016 (Resolução CNE/CES nº 5/2016), a observação do CC2020 é fundamental, pois ajuda a estruturar um currículo que integra hardware e software, diferenciando-o de cursos como Ciência da Computação e Engenharia de Software. Além disso, o alinhamento com as diretrizes do CC2020 assegura que os egressos tenham uma formação sólida, moderna e interdisciplinar, contemplando áreas emergentes como cibersegurança, computação paralela e distribuída, inteligência artificial e sistemas embarcados. Dessa forma, a adoção dessas diretrizes reforça a coerência curricular, moderniza o conteúdo didático, promove a internacionalização do ensino e amplia as oportunidades de empregabilidade dos alunos no mercado global da computação.

% \subsection{Competências em Engenharia de Computação}

% A área da computação abrange uma ampla gama de conhecimentos, que vão desde os impactos sociais da tecnologia até a implementação de sistemas complexos. O Quadro~\ref{tab:elementos_conhecimento} organiza esses elementos em categorias fundamentais, cobrindo aspectos que incluem segurança, modelagem de sistemas, infraestrutura, desenvolvimento de software e fundamentos de hardware. Esses tópicos são essenciais para a formação de profissionais que atuarão em diversos setores, garantindo que possuam habilidades técnicas sólidas e uma compreensão abrangente da área. Cada elemento descrito reflete tanto a teoria quanto as aplicações práticas necessárias para o desenvolvimento de soluções computacionais eficientes e inovadoras.

% % \newenvironment{ecptable}[3]{
% %     \renewcommand{\familydefault}{\sfdefault} \normalfont
% %     \begin{table}[ht]
% %     \centering
% %     \footnotesize
% %     \renewcommand{\arraystretch}{1.2} % Ajusta o espaçamento entre linhas
% %     \caption{#1}
% %     \label{#2}
% %     \begin{tabular}{|#3|}
% %     \hline
% % }{
% %     \hline
% %     \end{tabular}
% %     \renewcommand{\familydefault}{\rmdefault} \normalfont
% %     \end{table}
% % }

% \renewcommand{\familydefault}{\sfdefault} \normalfont{}
% \begin{table}[h!]
%     \centering
%     \caption{Elementos de Conhecimento em Computação \textsuperscript{\scriptsize\textcolor{AzulClaro}{CC2020, Tab.4.1}}}
%     \label{tab:elementos_conhecimento}
%     \tiny 
%     \renewcommand{\arraystretch}{1.2}
%     \begin{tabular}{|m{5.5cm}|m{8.5cm}|}
% % \begin{ecptable}{Elementos de Conhecimento em Computação (CC2020, tab.4.1)}{tab:elementos_conhecimento}{l|p{10cm}}
%     \hline
%     \rowcolor{AzulEscuro} \textcolor{white}{\textbf{Elemento de Conhecimento}} & \textcolor{white}{\textbf{Elaboração}} \\
%     \hline
%     \rowcolor{AzulClaro} \multicolumn{2}{|l|}{\textcolor{white}{\textbf{Usuários e Organizações}}} \\
%     \hline
%     $\cdot$ Questões Sociais e Prática Profissional & Impacto da tecnologia na sociedade, responsabilidades éticas dos profissionais de computação e conformidade com padrões de conduta profissional (ex: privacidade, sustentabilidade, inclusão). \\
%     \hline
%     $\cdot$ Política e Gestão de Segurança & Definição de políticas para proteção de dados, gestão de riscos e conformidade com regulamentações (ex: LGPD, ISO 27001). \\
%     \hline
%     $\cdot$ Gestão e Liderança de SI & Coordenação de equipes e recursos tecnológicos para alinhar sistemas de informação aos objetivos estratégicos da organização. \\
%     \hline
%     $\cdot$ Arquitetura Empresarial & Desenho e integração de sistemas, processos e infraestrutura para otimizar operações organizacionais. \\
%     \hline
%     $\cdot$ Gestão de Projetos & Aplicação de metodologias (ex: Agile, PMBOK) para planejar, executar e monitorar projetos de TI. \\
%     \hline
%     $\cdot$ Design de Experiência do Usuário (UX) & Criação de interfaces intuitivas e acessíveis, focadas nas necessidades e expectativas dos usuários finais. \\
%     \hline
%     \rowcolor{AzulClaro} \multicolumn{2}{|l|}{\textcolor{white}{\textbf{Modelagem de Sistemas}}} \\
%     \hline
%     $\cdot$ Questões e Princípios de Segurança & Identificação de vulnerabilidades e implementação de princípios de segurança (ex: confidencialidade, integridade) no ciclo de vida do sistema. \\
%     \hline
%     $\cdot$ Análise e Design de Sistemas & Tradução de requisitos em modelos funcionais e técnicos, utilizando técnicas como UML ou diagramas de fluxo. \\
%     \hline
%     $\cdot$ Análise de Requisitos e Especificações & Coleta, documentação e validação de necessidades de stakeholders para guiar o desenvolvimento. \\
%     \hline
%     $\cdot$ Gestão de Dados e Informação & Organização, armazenamento e recuperação de dados, incluindo uso de bancos de dados e governança de dados. \\
%     \hline
%     \rowcolor{AzulClaro} \multicolumn{2}{|l|}{\textcolor{white}{\textbf{Arquitetura e Infraestrutura de Sistemas}}} \\
%     \hline
%     $\cdot$ Sistemas e Serviços Virtuais & Implementação de ambientes virtualizados (ex: cloud computing, containers) para otimizar recursos físicos. \\
%     \hline
%     $\cdot$ Sistemas Inteligentes (IA) & Desenvolvimento de algoritmos e modelos de aprendizado de máquina para automação e tomada de decisão. \\
%     \hline
%     $\cdot$ Internet das Coisas (IoT) & Integração de dispositivos conectados à internet para coleta e processamento de dados em tempo real. \\
%     \hline
%     $\cdot$ Computação Paralela e Distribuída & Projeto de sistemas que executam tarefas simultaneamente em múltiplos processadores ou nós de rede. \\
%     \hline
%     $\cdot$ Redes de Computadores & Configuração e gerenciamento de protocolos, topologias e infraestrutura de comunicação (ex: TCP/IP, redes 5G). \\
%     \hline
%     $\cdot$ Sistemas Embarcados & Desenvolvimento de software e hardware integrados a dispositivos específicos (ex: sistemas automotivos, robótica). \\
%     \hline
%     $\cdot$ Tecnologia de Sistemas Integrados & Combinação de subsistemas heterogêneos (hardware/software) para funcionamento unificado. \\
%     \hline
%     $\cdot$ Tecnologias de Plataforma & Utilização de frameworks e ambientes pré-configurados para acelerar o desenvolvimento (ex: AWS, Kubernetes). \\
%     \hline
%     $\cdot$ Tecnologia e Implementação de Segurança & Aplicação de ferramentas e técnicas para proteger infraestrutura (ex: firewalls, criptografia). \\
%     \hline
%     \rowcolor{AzulClaro} \multicolumn{2}{|l|}{\textcolor{white}{\textbf{Desenvolvimento de Software}}} \\
%     \hline
%     $\cdot$ Qualidade de Software, Verificação e Validação & Garantia de que o software atende aos requisitos através de testes, revisões e métricas de qualidade (ex: testes unitários, CMMI). \\
%     \hline
%     $\cdot$ Processos de Software & Adoção de metodologias estruturadas para desenvolvimento (ex: DevOps, Scrum). \\
%     \hline
%     $\cdot$ Modelagem e Análise de Software & Criação de representações abstratas do sistema para identificar falhas e otimizar desempenho. \\
%     \hline
%     $\cdot$ Design de Software & Definição de arquitetura modular, padrões de projeto (ex: MVC) e boas práticas de codificação. \\
%     \hline
%     $\cdot$ Desenvolvimento Baseado em Plataformas & Utilização de ecossistemas pré-existentes (ex: Android Studio, Salesforce) para construção de aplicações. \\
%     \hline
%     \rowcolor{AzulClaro} \multicolumn{2}{|l|}{\textcolor{white}{\textbf{Fundamentos de Software}}} \\
%     \hline
%     $\cdot$ Gráficos e Visualização & Técnicas para representação visual de dados e criação de interfaces gráficas (ex: OpenGL, ferramentas de BI). \\
%     \hline
%     $\cdot$ Sistemas Operacionais & Estudo de gerenciamento de recursos (memória, processos) e interação entre hardware e software (ex: Linux, Windows). \\
%     \hline
%     $\cdot$ Estruturas de Dados, Algoritmos e Complexidade & Projeto e implementação de estruturas eficientes (ex: árvores, grafos, tabelas hash), desenvolvimento de algoritmos otimizados (ex: ordenação, busca, algoritmos em grafos) e/ou análise de eficiência computacional (ex: notação Big-O, escalabilidade, trade-offs entre tempo e espaço). \\
%     \hline
%     $\cdot$ Linguagens de Programação & Domínio de paradigmas (ex: orientação a objetos, funcional) e sintaxe de linguagens (ex: Python, Java). \\
%     \hline
%     $\cdot$ Fundamentos de Programação & Princípios básicos de lógica, sintaxe e depuração para construção de programas. \\
%     \hline
%     $\cdot$ Fundamentos de Sistemas Computacionais & Compreensão da interação entre hardware, software e redes em sistemas computacionais. \\
%     \hline
%     \rowcolor{AzulClaro} \multicolumn{2}{|l|}{\textcolor{white}{\textbf{Hardware}}} \\
%     \hline
%     $\cdot$ Arquitetura e Organização de Computadores & Projeto de componentes físicos (ex: CPU, memória) e sua interconexão para otimizar desempenho. \\
%     \hline
%     $\cdot$ Design Digital & Criação de circuitos lógicos e sistemas digitais usando portas lógicas e linguagens de descrição de hardware (ex: VHDL). \\
%     \hline
%     $\cdot$ Circuitos e Eletrônica & Estudo de componentes eletrônicos (ex: transistores, capacitores) e sua aplicação em sistemas computacionais. \\
%     \hline
%     $\cdot$ Processamento de Sinais & Técnicas para análise e manipulação de sinais analógicos/digitais (ex: filtros, compressão de áudio). \\
%     \hline
%     \end{tabular}
% \end{table}
% \renewcommand{\familydefault}{\rmdefault} \normalfont{}

%     Além do conhecimento técnico, a formação em computação exige competências fundamentais e profissionais que influenciam diretamente a atuação no mercado de trabalho. O Quadro~\ref{tab:elementos_fundamentais} destaca habilidades essenciais como pensamento crítico, colaboração, resolução de problemas e comunicação, bem como aspectos organizacionais como gestão de projetos e qualidade. Esses elementos são fundamentais para o sucesso de um profissional na área, permitindo que ele analise problemas complexos, trabalhe de maneira eficaz em equipe e se comunique com clareza em diferentes contextos. O desenvolvimento dessas habilidades complementa a formação técnica, tornando os profissionais mais preparados para enfrentar desafios multidisciplinares.

% \renewcommand{\familydefault}{\sfdefault} \normalfont{}
% \begin{table}[h!]
%     \centering
%     \caption{Elementos de Conhecimento Gerais \textsuperscript{\scriptsize\textcolor{AzulClaro}{CC2020, Tab.4.2}}}
%     \label{tab:elementos_fundamentais}
%     \tiny 
%     \renewcommand{\arraystretch}{1.2}
%     \begin{tabular}{|m{5cm}|m{9cm}|}
% \hline
% \rowcolor{AzulEscuro} \textcolor{white}{\textbf{Elemento de Conhecimento}} & \textcolor{white}{\textbf{Elaboração}} \\
% \hline
% $\cdot$ Pensamento Analítico e Crítico & Processo mental de simplificar informações complexas e tomar decisões \\
% \hline
% $\cdot$ Colaboração e Trabalho em Equipe & Dividir tarefas complexas e trabalhar em conjunto para completá-las \\
% \hline
% $\cdot$ Perspectivas Éticas e Interculturais & Visões éticas de diferentes pontos de vista em contextos humanos \\
% \hline
% $\cdot$ Matemática e Estatística & Uso abstrato de números e teorias, especialmente na análise de dados \\
% \hline
% $\cdot$ Priorização e Gestão de Múltiplas Tarefas & Organizar e priorizar múltiplas tarefas simultaneamente \\
% \hline
% $\cdot$ Comunicação Oral e Apresentação & Transmitir mensagens com eficácia usando recursos visuais e orais \\
% \hline
% $\cdot$ Resolução de Problemas e Solução de Erros & Buscar e resolver problemas de maneira lógica e organizada \\
% \hline
% $\cdot$ Organização e Planejamento de Projetos e Tarefas & Planejar e organizar projetos para atingir resultados \\
% \hline
% $\cdot$ Controle e Garantia de Qualidade & Usar métodos para identificar e prevenir defeitos \\
% \hline
% $\cdot$ Gestão de Relacionamentos & Manter o engajamento com clientes e parceiros \\
% \hline
% $\cdot$ Pesquisa e Aprendizagem Independente & Iniciar e conduzir projetos sem necessidade de direção \\
% \hline
% $\cdot$ Gestão de Tempo & Utilizar o tempo de maneira produtiva e eficiente \\
% \hline
% $\cdot$ Comunicação Escrita & Uso eficaz da escrita para comunicação entre pessoas e organizações \\
% \hline
% \end{tabular}
% \end{table}
% \renewcommand{\familydefault}{\rmdefault} \normalfont{}

% A aprendizagem em computação envolve diferentes níveis de complexidade cognitiva, que vão desde a memorização de conceitos até a criação de soluções inovadoras. A Taxonomia de Bloom organiza esses níveis de habilidade de forma hierárquica, permitindo estruturar o ensino de maneira progressiva. O Quadro~\ref{tab:habilidades_cognitivas} apresenta essas categorias, que incluem recordação, compreensão, aplicação, análise, avaliação e criação. Esse modelo é amplamente utilizado na educação para desenvolver materiais didáticos e avaliações que incentivem o aprendizado significativo, garantindo que os estudantes não apenas absorvam informações, mas também saibam utilizá-las para resolver problemas reais.

% \renewcommand{\familydefault}{\sfdefault} \normalfont{}
% \begin{table}[h!]
%     \centering
%     \caption{Níveis de Habilidade Cognitiva (Taxonomia de Bloom)  \textsuperscript{\scriptsize\textcolor{AzulClaro}{CC2020, Tab.4.3}}}
%     \label{tab:habilidades_cognitivas}
%     \tiny 
%     \renewcommand{\arraystretch}{1.2}
%     \begin{tabular}{|m{4cm}|m{10cm}|}
%     \hline
%     \rowcolor{AzulEscuro} \textcolor{white}{\textbf{Nível de Habilidade Cognitiva}} & \textcolor{white}{\textbf{Significado}} \\
%     \hline
%     $\cdot$ Recordação & Ato de relembrar materiais previamente aprendidos \\
%     \hline
%     $\cdot$ Compreensão & Capacidade de demonstrar entendimento (por meio de organização e comparação) \\
%     \hline
%     $\cdot$ Aplicação & Utilização do conhecimento para resolver problemas novos \\
%     \hline
%     $\cdot$ Análise & Exame e decomposição das informações em partes \\
%     \hline
%     $\cdot$ Avaliação & Julgamento sobre a validade e qualidade das informações \\
%     \hline
%     $\cdot$ Criação & Compilação das informações em padrões novos ou proposição de alternativas \\
%     \hline
%     \end{tabular}
%     \end{table}
%     \renewcommand{\familydefault}{\rmdefault} \normalfont{}

%     Além das habilidades técnicas e cognitivas, os profissionais da computação devem desenvolver disposições que orientam sua postura e atuação no ambiente profissional. O Quadro~\ref{tab:disposicoes_potenciais} apresenta características desejáveis, como adaptabilidade, colaboração, criatividade e responsabilidade, que impactam diretamente a eficiência e o sucesso no trabalho. Essas disposições são fundamentais para enfrentar os desafios do setor tecnológico, que exige constante inovação, trabalho em equipe e uma postura ética diante das transformações digitais. O cultivo dessas qualidades contribui para a formação de profissionais mais completos e preparados para lidar com os desafios dinâmicos da computação.

%     \renewcommand{\familydefault}{\sfdefault} \normalfont{}
% \begin{table}[h!]
%     \centering
%     \caption{Disposições Potenciais  \textsuperscript{\scriptsize\textcolor{AzulClaro}{CC2020, Tab.4.4}}}
%     \label{tab:disposicoes_potenciais}
%     \tiny 
%     \renewcommand{\arraystretch}{1.2}
%     \begin{tabular}{|m{4cm}|m{10cm}|}
%         \hline
%         \rowcolor{AzulEscuro} \textcolor{white}{\textbf{Disposição}} & \textcolor{white}{\textbf{Elaboração}} \\
%         \hline
%         $\cdot$ Adaptável & Flexível; ágil, ajusta-se em resposta à mudança \\
%         \hline
%         $\cdot$ Colaborativo & Colabora em equipe, demonstra disposição para trabalhar em conjunto \\
%         \hline
%         $\cdot$ Inventivo & Exploratório; busca soluções além do óbvio \\
%         \hline
%         $\cdot$ Meticuloso & Atento aos detalhes; minucioso, preciso \\
%         \hline
%         $\cdot$ Apaixonado & Convicção; forte comprometimento; envolvimento cativante \\
%         \hline
%         $\cdot$ Proativo & Atua com iniciativa; empreendedor; independente \\
%         \hline
%         $\cdot$ Profissional & Demonstra profissionalismo, discrição, comportamento ético e perspicácia \\
%         \hline
%         $\cdot$ Orientado por Propósito & Focado em metas; busca atingir objetivos; apresenta discernimento em termos de resultados e negócios \\
%         \hline
%         $\cdot$ Responsável & Exerce bom julgamento, usa discrição; age de maneira apropriada \\
%         \hline
%         $\cdot$ Solícito & Respeitoso; reage de forma rápida e positiva \\
%         \hline
%         $\cdot$ Autodirigido & Autônomo, motivado internamente, determinado, independente \\
%         \hline
%         \end{tabular}
%         \end{table}
%         \renewcommand{\familydefault}{\rmdefault} \normalfont{}


% pg. 108 pdf

% CE-CAE --- Circuits and Electronics
% CE-CAL — Computing Algorithms
% CE-CAO — Computer Architecture \& Organization
% CE-DIG — Digital Design
% CE-ESY — Embedded Systems
% CE-NWK — Computer Networks
% CE-PPP — Preparation for Professional Practice
% CE-SEC — Information Security
% CE-SGP — Signal Processing
% CE-SPE — Systems and Project Engineering
% CE-SRM — Systems Resource Management
% CE-SWD — Software Design

% \section{Conhecimentos em Engenharia de Computação}

% % pg. 65 pdf
% 1. Users and Organizations
% 1.1. Social Issues and Professional Practice 2 5
% --- 1.2. Security Policy and Management 1 3
% --- 1.3. IS Management and Leadership 0 2
% --- 1.4. Enterprise Architecture 0 1
% 1.5. Project Management 1 3
% --- 1.6. User Experience Design 1 3
% 2. Systems Modeling 
% --- 2.1. Security Issues and Principles 2 3
% --- 2.2. Systems Analysis \& Design 1 2
% 2.3. Requirements Analysis and Specification 1 2
% --- 2.4. Data and Information Management 1 2
% 3. Systems Architecture and Infrastructure
% 3.1. Virtual Systems and Services 1 3
% 3.2. Intelligent Systems (AI) 1 3
% *** 3.3. Internet of Things 2 4
% *** 3.4. Parallel and Distributed Computing 2 4
% 3.5. Computer Networks 2 4
% *** 3.6. Embedded Systems 3 5
% 3.7. Integrated Systems Technology 1 2
% --- 3.8. Platform Technologies 0 1
% 3.9. Security Technology and Implementation 2 3
% 4. Software Development
% 4.1. Software Quality, Verification and Validation 1 3
% 4.2. Software Process 1 2
% 4.3. Software Modeling and Analysis 1 3
% 4.4. Software Design 2 4
% 4.5. Platform-Based Development 0 2
% 5. Software Fundamentals
% 5.1. Graphics and Visualization 1 2
% 5.2. Operating Systems 2 4
% 5.3. Data Structures, Algorithms and Complexity 2 4
% 5.4. Programming Languages 2 3
% 5.5. Programming Fundamentals 2 4
% *** 5.6. Computing Systems Fundamentals 2 3
% 6. Hardware 
% *** 6.1. Architecture and Organization 4 5
% *** 6.2. Digital Design 4 5
% *** 6.3. Circuits and Electronics 4 5
% *** 6.4. Signal Processing 3 4 

No Capítulo~\ref{chp:apresentacao_acm}, foi apresentada a organização curricular de um curso de Engenharia de Computação alinhada com os elementos de conhecimento estabelecidos na CC2020. Para garantir uma formação abrangente, atualizada e globalizada, as disciplinas foram estruturadas de forma a cobrir os principais domínios da computação.
 
Nesta seção, são apresentados os quadros que relacionam cada categoria de conhecimento da computação aos componentes curriculares correspondentes do curso de \nomecursonovo, evidenciando como os conteúdos são distribuídos ao longo do curso.

Pela ordem de importância dentro da área de competências da Engenharia de Computação, temos:

\begin{enumerate}
\item Quadro~\ref{tab:cc2020_hardware}, Hardware (Tabela~\ref{tab:elementos_conhecimento}, item 6), com destaque para \textit{Arquitetura e Organização de Computadores}, \textit{Projetos Digitais}, \textit{Circuitos e Eletrônica} e \textit{Processamento de Sinais}.
\item Quadro~\ref{tab:cc2020_arquitetura_infraestrutura_sistemas}, Arquitetura e Infraestrutura de Sistemas (Tabela~\ref{tab:elementos_conhecimento}, item 3), com destaque \textit{Internet das Coisas (IoT)}, \textit{Computação Paralela e Distribuída} e \textit{Sistemas Embarcados}, além da inclusão de \textit{Inteligência Artificial} (ver Seção~\ref{sec:sobre_IA}).
\item Quadro~\ref{tab:cc2020_fundamentos_software}, Fundamentos de Software (Tabela~\ref{tab:elementos_conhecimento}, item 5), com destaque para \textit{Fundamentos de Sistemas Computacionais}.
\item Quadro~\ref{tab:cc2020_usuarios_organizacoes}, Usuários e Organizações (Tabela~\ref{tab:elementos_conhecimento}, item 1).
\item Quadro~\ref{tab:cc2020_modelagem_sistemas}, Modelagem de Sistemas (Tabela~\ref{tab:elementos_conhecimento}, item 2).
\item Quadro~\ref{tab:cc2020_desenvolvimento_software}, Desenvolvimento de Software (Tabela~\ref{tab:elementos_conhecimento}, item 4).
\end{enumerate}

Com relação à disciplinas obrigatórias do curso, eis os quadros com as classificações das mesmas:

\begin{longtblr}[
        theme = ecp,
        caption = {CC2020: Hardware},
        label = {tab:cc2020_hardware},]{
    colspec = {Q[l,m,wd=63mm]Q[l,m,wd=23mm]Q[c,m,wd=7mm]Q[c,m,wd=7mm]Q[c,m,wd=7mm]Q[c,m,wd=7mm]Q[c,m,wd=9mm]},
    rowhead = 1,
    row{odd} = {bg=CinzaClaro},
    row{1} = {bg=AzulEscuro, fg=white},
    row{13} = {bg=AzulClaro, fg=white},
    cells  = {font=\fontsize{10pt}{12pt}\selectfont},
    }
        \textbf{Componente} & \textbf{Código}  & \textbf{CH Teór.} & \textbf{CH Prát.}  & \textbf{CH EaD} & \textbf{CH Ext.} & \textbf{CH Total} \\
        Arquitetura e Organização de Computadores & FEELT!AOC & 60 & 0 & 0 & 0 & 60 \\
    Circuitos Elétricos I & FEELT31301 & 75 & 0 & 0 & 0 & 75 \\
    Experimental de Circuitos Elétricos I & FEELT31302 & 0 & 15 & 0 & 0 & 15 \\
    Elementos de Sistemas Computacionais & FEELT!ESC & 30 & 0 & 30 & 0 & 60 \\
    Eletrônica Analógica I & FEELT31401 & 60 & 0 & 0 & 0 & 60 \\
    Experimental de Eletrônica Analógica I & FEELT31402 & 0 & 30 & 0 & 0 & 30 \\
    Introdução à Prototipagem Eletrônica & FEELT!IPE & 0 & 30 & 0 & 0 & 30 \\
    Metrologia & FEELT31204 & 30 & 30 & 0 & 0 & 60 \\
    Processamento Digital de Sinais & FEELT31628 & 45 & 15 & 0 & 0 & 60 \\
    Sistemas Digitais & FEELT31409 & 30 & 0 & 0 & 0 & 30 \\
    Experimental de Sistemas Digitais & FEELT31410 & 0 & 30 & 0 & 0 & 30 \\*
    \textbf{TOTAL} &  & \textbf{330} & \textbf{150} & \textbf{30} & \textbf{0} & \textbf{510} \\
    \end{longtblr}
    

\begin{longtblr}[
        theme = ecp,
        caption = {CC2020: Arquitetura e Infraestrutura de Sistemas},
        label = {tab:cc2020_arquitetura_infraestrutura_sistemas},]{
    colspec = {Q[l,m,wd=63mm]Q[l,m,wd=23mm]Q[c,m,wd=7mm]Q[c,m,wd=7mm]Q[c,m,wd=7mm]Q[c,m,wd=7mm]Q[c,m,wd=9mm]},
    rowhead = 1,
    row{odd} = {bg=CinzaClaro},
    row{1} = {bg=AzulEscuro, fg=white},
    row{15} = {bg=AzulClaro, fg=white},
    cells  = {font=\fontsize{10pt}{12pt}\selectfont},
    }
        \textbf{Componente} & \textbf{Código}  & \textbf{CH Teór.} & \textbf{CH Prát.}  & \textbf{CH EaD} & \textbf{CH Ext.} & \textbf{CH Total} \\
        Aprendizagem Profunda e Modelos Generativos & FEELT!APMG & 30 & 15 & 0 & 0 & 45 \\
    Aprendizagem de Máquina & FEELT!AMAQ & 30 & 15 & 0 & 0 & 45 \\
    Arquitetura de Software Aplicada & FEELT!ASA & 15 & 45 & 0 & 0 & 60 \\
    Engenharia de Contexto em IA Generativa & FEELT!ECIA & 15 & 0 & 30 & 0 & 45 \\
    Fundamentos da Inteligência Artificial & FEELT!FIA & 30 & 15 & 0 & 0 & 45 \\
    IA Aplicada à Cibersegurança & FEELT!IACS & 15 & 0 & 30 & 0 & 45 \\
    Redes de Comunicações I & FEELT31724 & 45 & 15 & 0 & 0 & 60 \\
    Redes de Comunicações II & FEELT31831 & 45 & 15 & 0 & 0 & 60 \\
    Segurança de Sistemas Computacionais & FEELT!SEG & 15 & 0 & 30 & 0 & 45 \\
    Sistemas Computacionais em Tempo Real & FEELT!SCTR & 15 & 30 & 0 & 0 & 45 \\
    Sistemas Distribuídos & FEELT!SD & 45 & 0 & 0 & 0 & 45 \\
    Sistemas Embarcados I & FEELT31523 & 45 & 30 & 0 & 0 & 75 \\
    Sistemas e Controle & FEELT!SCON & 30 & 15 & 0 & 0 & 45 \\*
    \textbf{TOTAL} &  & \textbf{375} & \textbf{195} & \textbf{90} & \textbf{0} & \textbf{660} \\
    \end{longtblr}
    

\begin{longtblr}[
        theme = ecp,
        caption = {CC2020: Fundamentos de Software},
        label = {tab:cc2020_fundamentos_software},]{
    colspec = {Q[l,m,wd=63mm]Q[l,m,wd=23mm]Q[c,m,wd=7mm]Q[c,m,wd=7mm]Q[c,m,wd=7mm]Q[c,m,wd=7mm]Q[c,m,wd=9mm]},
    rowhead = 1,
    row{odd} = {bg=CinzaClaro},
    row{1} = {bg=AzulEscuro, fg=white},
    row{9} = {bg=AzulClaro, fg=white},
    cells  = {font=\fontsize{10pt}{12pt}\selectfont},
    }
        \textbf{Componente} & \textbf{Código}  & \textbf{CH Teór.} & \textbf{CH Prát.}  & \textbf{CH EaD} & \textbf{CH Ext.} & \textbf{CH Total} \\
        Análise de Algoritmos & FEELT!AA & 60 & 0 & 0 & 0 & 60 \\
    Estrutura de Dados & FEELT!ED & 60 & 0 & 0 & 0 & 60 \\
    Programação Funcional & FEELT!PF & 30 & 15 & 0 & 0 & 45 \\
    Programação Procedimental & FEELT!PP & 30 & 30 & 0 & 0 & 60 \\
    Programação Script & FEELT!PS & 30 & 30 & 0 & 0 & 60 \\
    Sistemas Operacionais & FEELT!SO & 45 & 0 & 0 & 0 & 45 \\
    Teoria da Computação & FEELT!TCOMP & 45 & 0 & 0 & 0 & 45 \\*
    \textbf{TOTAL} &  & \textbf{300} & \textbf{75} & \textbf{0} & \textbf{0} & \textbf{375} \\
    \end{longtblr}
    

\begin{longtblr}[
        theme = ecp,
        caption = {CC2020: Usuários e Organizações},
        label = {tab:cc2020_usuarios_organizacoes},]{
    colspec = {Q[l,m,wd=63mm]Q[l,m,wd=23mm]Q[c,m,wd=7mm]Q[c,m,wd=7mm]Q[c,m,wd=7mm]Q[c,m,wd=7mm]Q[c,m,wd=9mm]},
    rowhead = 1,
    row{odd} = {bg=CinzaClaro},
    row{1} = {bg=AzulEscuro, fg=white},
    row{6} = {bg=AzulClaro, fg=white},
    cells  = {font=\fontsize{10pt}{12pt}\selectfont},
    }
        \textbf{Componente} & \textbf{Código}  & \textbf{CH Teór.} & \textbf{CH Prát.}  & \textbf{CH EaD} & \textbf{CH Ext.} & \textbf{CH Total} \\
        Empreendedorismo e Inovação & FAGEN39401 & 30 & 0 & 0 & 0 & 30 \\
    Engenharia Econômica & IERI39601 & 30 & 0 & 0 & 0 & 30 \\
    Inteligência Artificial, Ética e Sociedade & FEELT!IAES & 30 & 0 & 0 & 0 & 30 \\
    Introdução à Engenharia de Computação & FEELT31103 & 30 & 0 & 0 & 0 & 30 \\*
    \textbf{TOTAL} &  & \textbf{120} & \textbf{0} & \textbf{0} & \textbf{0} & \textbf{120} \\
    \end{longtblr}
    

\begin{longtblr}[
        theme = ecp,
        caption = {CC2020: Modelagem de Sistemas},
        label = {tab:cc2020_modelagem_sistemas},]{
    colspec = {Q[l,m,wd=63mm]Q[l,m,wd=23mm]Q[c,m,wd=7mm]Q[c,m,wd=7mm]Q[c,m,wd=7mm]Q[c,m,wd=7mm]Q[c,m,wd=9mm]},
    rowhead = 1,
    row{odd} = {bg=CinzaClaro},
    row{1} = {bg=AzulEscuro, fg=white},
    row{4} = {bg=AzulClaro, fg=white},
    cells  = {font=\fontsize{10pt}{12pt}\selectfont},
    }
        \textbf{Componente} & \textbf{Código}  & \textbf{CH Teór.} & \textbf{CH Prát.}  & \textbf{CH EaD} & \textbf{CH Ext.} & \textbf{CH Total} \\
        Banco de Dados & FEELT!BD & 15 & 45 & 0 & 0 & 60 \\
    Programação Orientada a Objetos & FEELT!POO & 30 & 30 & 0 & 0 & 60 \\*
    \textbf{TOTAL} &  & \textbf{45} & \textbf{75} & \textbf{0} & \textbf{0} & \textbf{120} \\
    \end{longtblr}
    

\begin{longtblr}[
        theme = ecp,
        caption = {CC2020: Desenvolvimento de Software},
        label = {tab:cc2020_desenvolvimento_software},]{
    colspec = {Q[l,m,wd=63mm]Q[l,m,wd=23mm]Q[c,m,wd=7mm]Q[c,m,wd=7mm]Q[c,m,wd=7mm]Q[c,m,wd=7mm]Q[c,m,wd=9mm]},
    rowhead = 1,
    row{odd} = {bg=CinzaClaro},
    row{1} = {bg=AzulEscuro, fg=white},
    row{3} = {bg=AzulClaro, fg=white},
    cells  = {font=\fontsize{10pt}{12pt}\selectfont},
    }
        \textbf{Componente} & \textbf{Código}  & \textbf{CH Teór.} & \textbf{CH Prát.}  & \textbf{CH EaD} & \textbf{CH Ext.} & \textbf{CH Total} \\
        Engenharia de Software & FEELT!ESOF & 45 & 0 & 0 & 0 & 45 \\*
    \textbf{TOTAL} &  & \textbf{45} & \textbf{0} & \textbf{0} & \textbf{0} & \textbf{45} \\
    \end{longtblr}
    

\begin{longtblr}[
        theme = ecp,
        caption = {Carga Horária por Elemento de Conhecimento},%\protect\footnotemark[\value{footnote}]},
        label = {tab:ch_acm_semanal},
        note{$\ast$} = {\scriptsize A carga horária total de disciplinas de IA é de 255 horas.}
    ]{
    colspec = {Q[l,m,wd=80mm]Q[c,m,wd=30mm]Q[c,m,wd=30mm]},
    rowhead = 1,
    row{odd} = {bg=CinzaClaro},
    row{1} = {bg=AzulEscuro, fg=white},
    row{8} = {bg=gray, fg=white},
    row{13} = {bg=gray, fg=white},
    row{14} = {bg=AzulClaro, fg=white},
    cells  = {font = \fontsize{10pt}{12pt}\selectfont},
    } 
        \textbf{Conhecimento} & \textbf{Total (horas)} & \textbf{Percentual} \\
        Hardware & 510 & 14,8\% \\
    Arquitetura e Infraestrutura de Sistemas\TblrNote{$\ast$} & 660 & 19,1\% \\
    Fundamentos de Software & 375 & 10,9\% \\
    Usuários e Organizações & 120 & 3,5\% \\
    Modelagem de Sistemas & 120 & 3,5\% \\
    Desenvolvimento de Software & 45 & 1,3\% \\
   \textbf{Subtotal} &  \textbf{1830} & \textbf{53,0\%} \\    Formação Matemática e Física Aplicadas à Engenharia & 795 & 23,0\% \\
    Formação Optativa e Complementar & 180 & 5,2\% \\
    Atividades Curriculares de Extensão & 345 & 10,0\% \\
    Atividade de Conclusão de Curso & 300 & 8,7\% \\
   \textbf{Subtotal} &  \textbf{1620} & \textbf{47,0\%} \\   \textbf{TOTAL} &  \textbf{3450} & \textbf{100,0\%} \\
    \end{longtblr}

% \subsection*{Conhecimento em Hardware} Conhecimento em Hardware com dimensões:

% \listspace
% \begin{enumerate}[label=\Roman*]
%     \item \label{enum:hw-1} Arquitetura e Organização
%     \item \label{enum:hw-2} Design digital
%     \item \label{enum:hw-3} Circuitos e Eletrônica
%     \item \label{enum:hw-4} Processamento de Sinal
% \end{enumerate}

% \renewcommand{\familydefault}{\sfdefault} \normalfont
% \begin{longtblr}[
%     theme = ecp,
%     caption = {Conhecimento em Hardware},
%     label = {tab:conh_hardware},
% ]{
%   colspec = {Q[l,m,wd=65mm]Q[c,m,wd=24mm]Q[c,m,wd=10mm]Q[c,m,wd=20mm]},
%   rowhead = 1,
%   row{odd} = {bg=CinzaClaro},
%   row{1} = {bg=AzulEscuro, fg=white},
%   row{2} = {bg=AzulClaro, fg=white},
%   row{10} = {bg=AzulClaro, fg=white},
%   row{15} = {bg=AzulClaro, fg=white},
%   cells  = {font=\fontsize{10pt}{12pt}\selectfont},
% } 
%     % Componente & Código & CH Teór. & CH Prát.  & CH Rem. & CH Total \\
%     \textbf{Componente} & \textbf{Código}  & \textbf{CH Total} & \textbf{Dimensões}\\
%     \ \ \ \textit{Obrigatórias} &&&\\
%     Enriquecimento Instrumental & FEELT31109 & 30 & \ref{enum:hw-3} \\
%     Circuitos Eletro-Eletrônicos I & FEELT? & 60 & \ref{enum:hw-3}  \\
%     Circuitos Eletro-Eletrônicos I & FEELT? & 60 & \ref{enum:hw-3}  \\
%     Elementos de Sistemas Computacionais: Hardware & FEELT? & 45 & \ref{enum:hw-1},\ref{enum:hw-2}  \\
%     Arquitetura e Organização de Computadores & FEELT? & 60 & \ref{enum:hw-1}  \\
%     Sinais e Sistemas & FEELT32504 & 60 & \ref{enum:hw-4}  \\
%     Computação Embarcada & FEELT? & 75 & \ref{enum:hw-1},\ref{enum:hw-3}  \\
%     \ \ \ \textit{Optativas} &&&\\
%     Projetos Eletrônicos e Microprocessados & FEELT? & 45 & \ref{enum:hw-2},\ref{enum:hw-3} \\
%     Sistemas de Controle Realimentado & FEELT32603 & 60  & \ref{enum:hw-4} \\
%     Sistemas Computacionais em Tempo Real & FEELT31725 & 45  & \ref{enum:hw-2},\ref{enum:hw-4} \\
%     Sistemas Embarcados II & FEELT39015 & 60  & \ref{enum:hw-1},\ref{enum:hw-3} \\
%     \textbf{TOTAL} &  & \textbf{330} & 17,4\% (11,3\%|6,1\%) \\
% \end{longtblr}
% \renewcommand{\familydefault}{\rmdefault} \normalfont